\chapter{Introduction}

Le phytoplancton désigne l'ensemble des organismes unicellulaires photosynthétiques (« phyto ») dépourvus de capacité de déplacement autonome contre le courant (« plancton »). Ces organismes soutiennent les écosystèmes marins en constituant le premier maillon de la chaîne alimentaire, aussi appelée réseau trophique. Ils jouent par ailleurs un rôle fondamental dans le cycle du carbone, en structurant, par leur taille, leur composition et leur abondance, les apports de carbone vers les fonds marins. On parle de communauté phytoplanctonique pour désigner un assemblage particulier d'organismes phytoplanctoniques dans un espace et à un temps donné \cite{reynolds2006ecology}.

\vspace{0.3cm}

Le réchauffement climatique induit des changements rapides dans l'océan Austral : l'eau se réchauffe et s'adoucit sous l'effet de la fonte des glaces, les champs de courants et de vents se modifient, la stratification de la colonne d'eau s'accentue, etc. \cite{sallee2018} \cite{abram2025Antarctic}. Ces changements ne sont pas sans conséquences pour les communautés phytoplanctoniques \cite{deppeler2017southern}. Des études récentes montrent des changements de phénologie \cite{thomalla2023Widespread}\cite{boyce2017environmental}, d'abondance \cite{hayward2024Antarctic} et de structuration des communautés phytoplanctoniques \cite{adroli2026}.

\vspace{0.3cm}

Dans l'océan Austral, ces communautés sont dominées par les diatomées, larges cellules phytoplanctoniques (microphytoplancton, 20-200 $\mu m$) à forte valeur énergétique \cite{treguer1992} \cite{heidemann2024drivers}. Les modifications induites par le réchauffement climatique tendraient à augmenter la biomasse phytoplanctonique \cite{delcastillo2019greener} \cite{adroli2026} tout en diminuant la contribution relative des diatomées au profit du nano- et picophytoplancton, organismes à plus faible valeur énergétique \cite{krumhardt2022climate}\cite{sari2025influence}. Les diatomées soutiennent un réseau trophique court et hautement productif \cite{treguer1992}. Elles sont consommées par le grand zooplancton (copépodes, krill) et le micronecton (petits poissons), regroupés sous l'appellation de Niveaux Trophiques Intermédiaires (NTI). Ces organismes sont eux-mêmes consommés par de plus gros poissons pélagiques, à leur tour proies de plus grands prédateurs. Un changement de structuration des communautés phytoplanctoniques vers une plus forte proportion de petit phytoplancton favoriserait ainsi le développement d'un zooplancton de plus petite taille, lui-même consommé par des organismes plus petits \cite{krumhardt2022climate}. Ainsi, on observerait un allongement du réseau trophique et donc une modification du transfert d'énergie à chacun des ses maillons jusqu'aux prédateurs supérieurs \cite{mestre2020decadal} : on parle de cascade trophique. Un changement d'assemblage des communautés phytoplanctonique impacterait donc l'ensemble de l'écosystème.


\vspace{0.3cm}

Caractériser ces changements est par conséquent nécessaire, mais difficile dans les écosystèmes pélagiques \footnote{Ensemble des organismes et de leur environnement vivant dans la colonne d'eau du large, c'est-à-dire en pleine eau, loin du fond et des côtes, par opposition au domaine benthique lié au fond marin.} : ces milieux étant très vastes et peu accessibles, leur échantillonnage à grande échelle est logistiquement couteux. L'acoustique active constitue un moyen d'étude des NTI dans ces zones difficiles d'accès. Elle consiste à émettre une onde, à la réceptionner, puis à analyser le signal reçu (principe du sonar) \cite{benoitbird2016ecological}. L'écho reçu résulte de la réflexion du son émis sur un objet (le fond marin, par exemple), un organisme, ou encore une discontinuité physique dans la colonne d'eau (changement de densité). La connaissance de la vitesse locale de propagation du son dans l'eau permet, à partir du délai de réception du signal et de son intensité, de renseigner la profondeur, l'abondance et d'inférer sur la composition taxonomique des organismes présents \cite{benoitbird2016ecological}. Le Nautical Area Scattering Coefficient (NASC), calculé à partir des échos acoustiques reçus, permet de résumer la rétrodiffusion acoustique sur un ensemble de profondeur considérées en un certain point de l'espace \cite{maclennan2002}. Il sert donc de proxy de la biomasse de NTI à cette localisation. L'océan Austral étant particulièrement vaste et difficile d'accès, les campagnes d'échantillonnages restent limitées et, dans ce contexte, il est particulièrement pertinent de développer des modèles prédictifs permettant de généraliser les estimations faites à partir des données déjà acquises.

\vspace{0.3cm}

Cette étude à cherché à développer un modèle prédictif de la biomasse des NTI en fonction de l'assemblage des communautés phytoplanctoniques et de l'activité de mésoéchelle. Plus précisément, nous avons modélisé le NASC avec des Forêts Aléatoires en fonction de 9 variables de concentration pigmentaires ($P_1, P_2, \ldots, P_9$), 2 variables physique, la température et salinité (FOD), et 1 variable capturant l'hydrodynamisme (FTLE). Ainsi, nous avons cherché à documenter la relation $NASC \sim FTLE + FOD + P_1 + P_2 + \ldots + P_9$. Ce modèle, développé sur le secteur indien de l'océan Austral à partir de données d'acoustique active collectées lors des campagnes OBSAUSTRAL 2018, 2021, 2022 et 2023, à permis de créer des cartes prédictives de NASC sur l'ensemble de la région et de mettre en évidence les variables environnementales clés pour la prédiction. L'objectif étant d'informer les conséquences de la restructuration des communautés phytoplanctoniques sur les NTI du secteur Indien de l'océan Austral.

\vspace{0.3cm}

La section \textbf{Matériel} détaillera la nature des données utilisées (acoustique active, température, salinité, champs de courants et concentrations pigmentaires). La section \textbf{Méthode} présentera le calcul des variables d'intérêt NASC, FTLE et FOD. Cette section décrira ensuite brièvement la méthode Random Forest, le calcul de l'importance des variables et le pipeline d'apprentissage mis en place. La section \textbf{Résultat} présentera (1) les cartes de NASC prédites (2) les variables environnementales identifiées comme importantes pour la prédiction du NASC (3) l'interprétation du modèle. La \textbf{Discussion }portera sur la fiabilité du modèle, la difficultée liée à la nature même des données et la mise en perspective avec la littérature de l'effet structurant des variables environnementales considérées et de la nature des communautés phytoplanctoniques.
